\documentclass[11pt,a4paper]{article}
\usepackage{amsmath,amssymb,amsthm}
\usepackage[margin=25mm]{geometry}
\usepackage{xeCJK}
\setCJKmainfont{Hiragino Mincho ProN}
\setCJKsansfont{Hiragino Sans}
\usepackage[hidelinks]{hyperref}
\usepackage[nameinlink,noabbrev]{cleveref}
\renewcommand{\proofname}{証明}
\newtheorem{definition}{定義}[section]
\newtheorem{claim}[definition]{主張}
\newtheorem{theorem}[definition]{定理}
\theoremstyle{remark}
\newtheorem{remark}[definition]{注意}
\newtheorem{structurednote}[definition]{ノート}
\crefname{definition}{定義}{定義}
\crefname{claim}{主張}{主張}
\crefname{theorem}{定理}{定理}
\title{有限双曲曲面上の Ising 模型の可算構造}
\author{}
\date{}
\begin{document}
\maketitle
\part{有限グラフ上の Ising 多項式}

\begin{definition}[有限グラフの入力]\label{lab:def_finite_graph_input}
有限グラフの入力を、空でない有限集合 
$V$
、有限集合 
$E$
、二つの写像 
$\partial_0,\partial_1:E\to V$
 の組 
$G=(V,E,\partial_0,\partial_1)$
 と定める。任意の 
$e\in E$
 について 
$\partial_0(e)\ne\partial_1(e)$
 を仮定する。異なる辺が同じ二頂点を結ぶことは許す。

\end{definition}

\begin{definition}[スピン配位集合]\label{lab:def_spin_configuration_set}
有限グラフ 
$G$
 のスピン配位集合を

\[\mathcal S_G:=\{\sigma\mid \sigma:V\to\{-1,+1\}\ \text{は写像}\}\]
と定める。
$\mathcal S_G$
 は有限集合であり、
$|\mathcal S_G|=2^{|V|}$
 である。ここで 
$|X|\in\mathbb N$
 は有限集合 
$X$
 の元の個数を表す。

\end{definition}

\begin{definition}[破れ辺集合と破れ辺数]\label{lab:def_broken_edge_set}
配位 
$\sigma\in\mathcal S_G$
 の破れ辺集合と破れ辺数を

\[B_G(\sigma):=\{e\in E\mid \sigma(\partial_0(e))\ne\sigma(\partial_1(e))\},\qquad b_G(\sigma):=|B_G(\sigma)|\in\mathbb N\]
と定める。端点写像と配位は 
\cref{lab:def_finite_graph_input}
 と 
\cref{lab:def_spin_configuration_set}
 で定めた。

\end{definition}

\begin{definition}[破れ辺数の多重度]\label{lab:def_broken_edge_multiplicity}
各 
$m\in\{0,1,\ldots,|E|\}$
 に対し多重度を

\[\Omega_G(m):=|\{\sigma\in\mathcal S_G\mid b_G(\sigma)=m\}|\in\mathbb N\]
と定める。破れ辺数は 
\cref{lab:def_broken_edge_set}
 で定めた。

\end{definition}

\begin{definition}[Ising 分配多項式]\label{lab:def_ising_partition_polynomial}
不定元 
$x$
 に対し Ising 分配多項式を

\[Z_G(x):=\sum_{\sigma\in\mathcal S_G}x^{b_G(\sigma)}\in\mathbb Z[x]\]
と定める。これは多項式そのものであり、数を代入した値とは区別する。

\end{definition}

\begin{claim}[多重度による係数表示]\label{lab:claim_partition_polynomial_coefficient_expansion}
\[Z_G(x)=\sum_{m=0}^{|E|}\Omega_G(m)x^m\in\mathbb Z[x].\]
\end{claim}
\begin{proof}
\[\begin{aligned}
Z_G(x)
&=\sum_{\sigma\in\mathcal S_G}x^{b_G(\sigma)}
&&\bigl(\because\ \text{Ising 分配多項式の定義}\bigr)\\
&=\sum_{m=0}^{|E|}\ \sum_{\substack{\sigma\in\mathcal S_G\\ b_G(\sigma)=m}}x^m
&&\bigl(\because\ \mathcal S_G\text{ を }b_G\text{ の値ごとのファイバーへ分割}\bigr)\\
&=\sum_{m=0}^{|E|}\Omega_G(m)x^m
&&\bigl(\because\ \text{多重度の定義}\bigr).
\end{aligned}\]
最初の等号は 
\cref{lab:def_ising_partition_polynomial}
、最後の等号は 
\cref{lab:def_broken_edge_multiplicity}
 を用いた。

\end{proof}

\begin{claim}[係数総和]\label{lab:claim_partition_polynomial_value_at_one}
\[Z_G(1)=\sum_{m=0}^{|E|}\Omega_G(m)=2^{|V|}.\]
\end{claim}
\begin{proof}
\[\begin{aligned}
Z_G(1)
&=\sum_{\sigma\in\mathcal S_G}1^{b_G(\sigma)}
&&\bigl(\because\ \text{Ising 分配多項式の定義}\bigr)\\
&=\sum_{\sigma\in\mathcal S_G}1
&&\bigl(\because\ 1^n=1\ \text{for every }n\in\mathbb N\bigr)\\
&=|\mathcal S_G|
&&\bigl(\because\ \text{有限集合の元の個数の定義}\bigr)\\
&=2^{|V|}
&&\bigl(\because\ \text{各頂点で二つのスピン値を独立に選ぶ}\bigr).
\end{aligned}\]
多重度の和との等号は 
\cref{lab:claim_partition_polynomial_coefficient_expansion}
 に 
$x=1$
 を代入して得る。

\end{proof}

\part{形式的高温展開}

\begin{definition}[辺のスピン符号]\label{lab:def_edge_spin_sign}
各 
$\sigma\in\mathcal S_G$
 と 
$e\in E$
 に対し

\[s_G(\sigma,e):=\sigma(\partial_0(e))\sigma(\partial_1(e))\in\{-1,+1\}\]
と定める。

\end{definition}

\begin{definition}[形式的辺重み和]\label{lab:def_formal_edge_weight_sum}
独立な不定元 
$u,v$
 に対し

\[H_G(u,v):=\sum_{\sigma\in\mathcal S_G}\prod_{e\in E}\bigl(u+v\,s_G(\sigma,e)\bigr)\in\mathbb Z[u,v]\]
と定める。指数関数、双曲関数、実温度はこの定義に含まれない。

\end{definition}

\begin{definition}[辺部分集合の境界偶奇]\label{lab:def_mod_two_boundary_parity}
辺部分集合 
$A\subseteq E$
 と頂点 
$w\in V$
 に対し

\[\partial_G(A)(w):=\left|\{e\in A\mid \partial_0(e)=w\ \text{または}\ \partial_1(e)=w\}\right|\bmod 2\in\mathbb F_2\]
と定める。各辺の二端点は異なるので、一つの辺を同じ頂点で二重に数えない。

\end{definition}

\begin{definition}[偶辺部分集合]\label{lab:def_even_edge_subset}
\[\mathcal Z_1(G):=\{A\subseteq E\mid \partial_G(A)(w)=0\ \text{for every }w\in V\}\]
と定める。これは有限集合である。境界偶奇は 
\cref{lab:def_mod_two_boundary_parity}
 で定めた。

\end{definition}

\begin{definition}[偶部分グラフ多項式]\label{lab:def_even_subgraph_polynomial}
\[Q_G(u,v):=\sum_{A\in\mathcal Z_1(G)}u^{|E|-|A|}v^{|A|}\in\mathbb Z[u,v]\]
と定める。

\end{definition}

\begin{theorem}[形式的高温展開の有限和恒等式]\label{lab:theorem_formal_high_temperature_expansion}
\[H_G(u,v)=2^{|V|}Q_G(u,v)\in\mathbb Z[u,v].\]
\end{theorem}
\begin{proof}
積を辺部分集合ごとに展開すると

\[\begin{aligned}
H_G(u,v)
&=\sum_{\sigma\in\mathcal S_G}\ \sum_{A\subseteq E}u^{|E|-|A|}v^{|A|}\prod_{e\in A}s_G(\sigma,e)
&&\bigl(\because\ \text{有限積に分配則を適用}\bigr)\\
&=\sum_{A\subseteq E}u^{|E|-|A|}v^{|A|}\sum_{\sigma\in\mathcal S_G}\prod_{e\in A}s_G(\sigma,e)
&&\bigl(\because\ \text{二つの有限和の順序交換}\bigr).
\end{aligned}\]
$A\in\mathcal Z_1(G)$
 なら、各頂点のスピンは積に偶数回現れるので内側の積は全ての配位で 
$1$
 となり、内側の和は 
$2^{|V|}$
 である。

$A\notin\mathcal Z_1(G)$
 なら、境界偶奇が 
$1$
 である頂点 
$w$
 が存在する。配位の 
$w$
 での値だけを反転する写像は 
$\mathcal S_G$
 上の不動点を持たない対合であり、対になった二項の積の符号が逆になるので、内側の和は 
$0$
 である。

\[\begin{aligned}
H_G(u,v)
&=\sum_{A\in\mathcal Z_1(G)}u^{|E|-|A|}v^{|A|}\,2^{|V|}
&&\bigl(\because\ \text{上の偶奇による二場合}\bigr)\\
&=2^{|V|}Q_G(u,v)
&&\bigl(\because\ \text{偶部分グラフ多項式の定義}\bigr).
\end{aligned}\]
用いた定義は 
\cref{lab:def_formal_edge_weight_sum}
、
\cref{lab:def_edge_spin_sign}
、
\cref{lab:def_even_edge_subset}
、
\cref{lab:def_even_subgraph_polynomial}
 である。

\end{proof}
\end{document}
